\section{Problem and Motivation}
\label{sec:intro}

Image captioning models write fluent sentences, and fluency hides their mistakes. A dog carrying a frisbee is described as ``a dog playing with a ball''; a park scene gains a bench that is not in the picture. Rohrbach \emph{et al.}~\cite{rohrbach2018object} showed that such \emph{object hallucinations} follow the statistics of the training captions more than the pixels, and that a model can improve on sentence-level metrics such as CIDEr while hallucinating more. For a reader who cannot see the image, for instance a screen-reader user, a wrongly named object is worse than a missing one, because nothing in the sentence signals that it should be doubted.

\textbf{Problem.} We work at the level of \emph{object claims}. A claim is a concrete noun in a caption that belongs to a fixed vocabulary $V$ (Section~\ref{sec:design-claims}). All person words (\emph{man}, \emph{girl}, \emph{people}, \dots) form a single class \emph{person} that is checked for presence only; gender and age are never verified or scored, because detector-based checks of them are unreliable and raise fairness concerns~\cite{hendricks2018women}. A claim is \emph{incorrect} if no such object is visible or if another object has been given its name. Given an image and a draft caption, we want a new caption with fewer incorrect claims that keeps the correct ones and does not lose caption quality.

\textbf{Approach.} \emph{Look--Check--Retell} (LCR) is a two-pass captioner (Fig.~\ref{fig:pipeline}). A frozen BLIP model~\cite{li2022blip} drafts a caption. \emph{Look} uses BLIP's own decoder cross-attention to pick an image region for each claim. \emph{Check} asks a separate frozen open-vocabulary detector, OWLv2~\cite{minderer2023scaling}, what that region contains and writes a small \emph{verification record}. \emph{Retell} regenerates the caption with BLIP, which now also attends to the records through a small trainable gated pathway. The detector is itself imperfect and its scores are not calibrated probabilities~\cite{guo2017calibration}, so the pathway has to learn \emph{when} evidence should change the caption; it must not copy the detector blindly.

\textbf{Research question.} Can a captioner learn to use noisy, claim-specific verification evidence to improve its captions while rejecting misleading evidence, when the checked regions are chosen by its own attention?

\textbf{Scope of this report.} The proposal fixed seven hypotheses (H1--H6 and H1b) and a 12-week plan. This report covers the first part of that plan. We describe the design decisions taken while turning the proposal into a working system and the reasons behind them (Section~\ref{sec:design}), and we report what has been measured so far (Section~\ref{sec:experiments}). All experiments use the Flickr8k training and validation splits only; the test split has not been touched, as the protocol requires. Several of our measurements use the reference captions as a stand-in for human judgment. We label every such number as a proxy, and we list what remains open in Section~\ref{sec:conclusion}.
